% ============================================================
%  PETRIFICATION PROJECT — LaTeX Article Template
%  θ-transform (Krasnoselskii-Mann) framework
%
%  Journal targets
%  ---------------
%  Short proofs / letters (≤4 pages):
%    Physical Review Letters (PRL)
%    → \documentclass[aps,prl,twocolumn,superscriptaddress,nofootinbib]{revtex4-2}
%
%  Research articles (≥6 pages):
%    Physical Review E (nonlinear dynamics / chaos / comp. physics)
%    → \documentclass[aps,pre,twocolumn,superscriptaddress,nofootinbib,longbibliography]{revtex4-2}
%
%  Computational physics focus:
%    Computer Physics Communications (CPC)
%    → switch to elsarticle class (see comment block at end)
%
%  Build
%  -----
%    pdflatex template
%    bibtex   template
%    pdflatex template
%    pdflatex template
%  or simply: make (uses Makefile in this directory)
% ============================================================

\documentclass[aps,pre,twocolumn,superscriptaddress,nofootinbib,longbibliography]{revtex4-2}

% ---- Unicode declarations (pdflatex + UTF-8 source) --------
% Declare Greek letters and math symbols used in text mode.
\DeclareUnicodeCharacter{03B8}{\ensuremath{\theta}}   % θ
\DeclareUnicodeCharacter{039B}{\ensuremath{\Lambda}}  % Λ
\DeclareUnicodeCharacter{03B5}{\ensuremath{\varepsilon}} % ε
\DeclareUnicodeCharacter{2264}{\ensuremath{\leq}}     % ≤
\DeclareUnicodeCharacter{2265}{\ensuremath{\geq}}     % ≥
\DeclareUnicodeCharacter{2248}{\ensuremath{\approx}}  % ≈
\DeclareUnicodeCharacter{2192}{\ensuremath{\to}}      % →
\DeclareUnicodeCharacter{00B7}{\ensuremath{\cdot}}    % ·

% ---- Core packages -----------------------------------------
\usepackage{amsmath}
\usepackage{amssymb}
\usepackage{amsfonts}
\usepackage{amsthm}          % theorem environments
\usepackage{graphicx}        % \includegraphics
\usepackage{dcolumn}         % aligned decimal columns in tables
\usepackage{bm}              % bold math symbols (\bm)
\usepackage{xcolor}
\usepackage{booktabs}        % \toprule / \midrule / \bottomrule
\usepackage{multirow}        % multirow cells in tables
\usepackage{siunitx}         % SI units (\SI{}, \num{})
\usepackage{hyperref}        % clickable refs — load last
\hypersetup{
  colorlinks = true,
  linkcolor  = blue!60!black,
  citecolor  = green!50!black,
  urlcolor   = blue!60!black,
}

% ============================================================
%  NOTATION MACROS
%
%  The Krasnoselskii-Mann relaxation parameter is denoted θ
%  (theta) throughout, consistent with the K-M iteration
%  literature.  The position-dependent generalisation θ(x) is
%  the "gauged" version introduced in this work.
%
%  NOTE ON NAMING CONFLICT
%  In the orbit-representation results (Sec. on closed-form
%  logistic orbits) the angle variable x = sin²(π φ) uses φ
%  (phi) to avoid collision with the K-M parameter θ.
% ============================================================

% --- K-M parameter and transform ---
\newcommand{\param}{\theta}                        % K-M relaxation parameter
\newcommand{\paramx}{\theta(x)}                    % position-dependent form
\newcommand{\optparam}{\theta^{*}}                 % optimal value
\newcommand{\transform}{g_{\param}}                % transformed map (subscript)
\newcommand{\transformof}[1]{g_{\param}(#1)}       % g_θ(·)
\newcommand{\KMdef}{%                              % inline definition
  g_{\param}(x) \;=\; \param\,f(x) + (1-\param)\,x}

% --- Fixed points ---
\newcommand{\FP}{x^{*}}
\newcommand{\FPi}[1]{x^{*}_{#1}}                  % indexed: \FPi{0}, \FPi{1}

% --- Dynamical quantities ---
\newcommand{\Lyap}{\Lambda}                        % Lyapunov exponent
\newcommand{\Lyapof}[1]{\Lambda\!\left(#1\right)}  % Λ(θ)
\newcommand{\orbit}[1]{\mathcal{O}(#1)}            % orbit of a point

% --- Perturbation / potential notation ---
\newcommand{\fzero}{f_0}                           % unperturbed map
\newcommand{\feps}{f_{\varepsilon}}                % perturbed map
\newcommand{\Vmodel}{V_{\mathrm{model}}}
\newcommand{\Vpert}{V_{\mathrm{pert}}}
\newcommand{\Vtrue}{V_{\mathrm{true}}}

% --- Transfer / Perron-Frobenius operator ---
\newcommand{\PF}{\mathcal{L}}                      % L_f
\newcommand{\PFparam}{\mathcal{L}_{g_{\param}}}    % L_{g_θ}

% --- Convenient shorthands ---
\newcommand{\KM}{Krasnoselskii--Mann}              % hyphenated name
\newcommand{\abs}[1]{\left|#1\right|}
\newcommand{\norm}[1]{\left\|#1\right\|}
\newcommand{\bigO}[1]{\mathcal{O}\!\left(#1\right)}
\newcommand{\reals}{\mathbb{R}}
\newcommand{\naturals}{\mathbb{N}}
\newcommand{\integers}{\mathbb{Z}}
\newcommand{\hilbert}{\mathcal{H}}

% --- Matrix / operator notation (for Dyson / KB papers) ---
\newcommand{\Sig}{\Sigma}                          % self-energy
\newcommand{\Green}{G}                             % Green's function
\newcommand{\Greenz}{G_0}                          % bare Green's function

% --- Plasma / fusion notation (activate for fusion papers) ---
%\newcommand{\vth}{v_{\mathrm{th}}}
%\newcommand{\wpe}{\omega_{\mathrm{pe}}}
%\newcommand{\kB}{k_{\mathrm{B}}}
%\newcommand{\betaPlasma}{\beta_{\mathrm{p}}}

% ============================================================
%  THEOREM ENVIRONMENTS  (amsthm)
% ============================================================

\theoremstyle{plain}
\newtheorem{theorem}{Theorem}[section]
\newtheorem{proposition}[theorem]{Proposition}
\newtheorem{corollary}[theorem]{Corollary}
\newtheorem{lemma}[theorem]{Lemma}

\theoremstyle{definition}
\newtheorem{definition}[theorem]{Definition}

\theoremstyle{remark}
\newtheorem{remark}[theorem]{Remark}

% ============================================================
%  DOCUMENT BEGINS
% ============================================================

\begin{document}

% ------------------------------------------------------------
%  FRONT MATTER
% ------------------------------------------------------------

\title{[TITLE: concise, keyword-rich, ≤12 words]}

\author{Keenan M.~Stone}
\affiliation{[AFFILIATION]}
\email{[EMAIL]}

% Add co-authors if applicable:
% \author{Second Author}
% \affiliation{[AFFILIATION 2]}

\date{\today}

\begin{abstract}
% Target: 150–250 words (PRL: ≤200 words).
% Structure: (1) context/problem, (2) what this paper does,
%            (3) key result in one sentence, (4) significance.
%
% Example opening: "We study the [problem] using the
% θ-transform $\KMdef$, a member of the
% \KM{} (KM) iteration family~\cite{Krasnoselskii1955,Mann1953}.
% ..."
[ABSTRACT PLACEHOLDER]
\end{abstract}

\keywords{Krasnoselskii-Mann iteration, fixed-point dynamics,
          θ-transform, logistic map, chaos control,
          relaxation methods, nonlinear dynamics}

\maketitle

% ============================================================
\section{Introduction}
\label{sec:intro}
% ============================================================

% Paragraph 1: broad context
% Paragraph 2: the specific open problem
% Paragraph 3: what this paper does (2–3 sentences)
% Paragraph 4: outline

The \KM{} (KM) iteration~\cite{Krasnoselskii1955,Mann1953} is a
fundamental relaxation scheme for approximating fixed points of
nonexpansive operators.  Given an operator $T$ and an initial
guess $x_0$, the sequence
%
\begin{equation}
  x_{n+1} = (1-\param_n)\,x_n + \param_n\,T(x_n),
  \label{eq:KM-iteration}
\end{equation}
%
with $\param_n \in (0,1)$, converges to a fixed point of $T$
under mild conditions~\cite{Berinde2007}.  For a smooth
self-map $f : I \to I$ with constant parameter $\param$,
Eq.~\eqref{eq:KM-iteration} defines the \emph{θ-transform}
%
\begin{equation}
  \KMdef,
  \label{eq:theta-transform}
\end{equation}
%
which preserves every fixed point of $f$ while rescaling their
local stability (Sec.~\ref{sec:theory}).
This paper investigates [STATE FOCUS].

The paper is organised as follows.
Section~\ref{sec:theory} states the core fixed-point
preservation and stability propositions.
Section~\ref{sec:results} presents [RESULTS].
Section~\ref{sec:discussion} discusses [DISCUSSION].
Section~\ref{sec:conclusions} concludes.

% ============================================================
\section{The θ-Transform: Definitions and Properties}
\label{sec:theory}
% ============================================================

\subsection{Fixed-point preservation and stability scaling}
\label{ssec:basic}

\begin{definition}[θ-transform]
\label{def:theta-transform}
Let $f : I \to I$ be a smooth self-map of an interval
$I \subset \reals$.  For $\param \in \reals$, the
\emph{θ-transform} of $f$ with parameter $\param$ is
$\KMdef$.
\end{definition}

\begin{proposition}[Fixed-point preservation]
\label{prop:fp-preservation}
For every $\param \in \reals$, $\transform$ and $f$ share the
same fixed-point set: $\transformof{\FP} = \FP$ if and only if
$f(\FP) = \FP$.
\end{proposition}
%
\begin{proof}
$\transformof{\FP} - \FP = \param(f(\FP) - \FP)$.  For
$\param \neq 0$ this vanishes iff $f(\FP) = \FP$; for
$\param = 0$ the transform collapses to the identity and
every point is trivially fixed.
\end{proof}

\begin{proposition}[Stability scaling]
\label{prop:stability}
At a fixed point $\FP$ of $f$,
%
\begin{equation}
  g_\param'(\FP) \;=\; 1 - \param\bigl(1 - f'(\FP)\bigr).
  \label{eq:stability-scaling}
\end{equation}
%
In particular, the derivative vanishes — giving maximally
attracting convergence — at the \emph{optimal parameter}
%
\begin{equation}
  \optparam = \frac{1}{1 - f'(\FP)}.
  \label{eq:optimal-theta}
\end{equation}
\end{proposition}
%
\begin{proof}
Differentiate $g_\param(x) = \param f(x) + (1-\param)x$:
$g_\param'(x) = \param f'(x) + (1-\param)$.  Evaluate at $x = \FP$
and simplify.
\end{proof}

\begin{corollary}[Stability inversion]
\label{cor:inversion}
A fixed point $\FP$ that is unstable under $f$
(i.e., $|f'(\FP)| > 1$) is stabilised under $\transform$
for $\param \in (0, \optparam)$ when $f'(\FP) > 1$, and for
$\param \in (\optparam, 0)$ when $f'(\FP) < -1$.
\end{corollary}

\subsection{Position-dependent θ(x)}
\label{ssec:position-dep}

Making $\param$ position-dependent ($\param \to \paramx$) is the
analogue of gauging a global symmetry: the composition rule
$\Gamma_\param \Gamma_\beta = \Gamma_{\param\beta}$ holds
pointwise.  The derivative now contains an additional
\emph{connection term}:
%
\begin{equation}
  g'(x) = \theta'(x)\bigl(f(x) - x\bigr)
           + 1 + \paramx\bigl(f'(x) - 1\bigr).
  \label{eq:position-dep-derivative}
\end{equation}
%
\begin{proposition}[Gauge invariance of fixed points]
\label{prop:gauge-invariance}
The connection term $\theta'(x)(f(x)-x)$ in
Eq.~\eqref{eq:position-dep-derivative} vanishes at every fixed
point, so the stability formula Eq.~\eqref{eq:stability-scaling}
holds locally regardless of the profile $\paramx$.
\end{proposition}

% ============================================================
\section{[RESULTS — PAPER-SPECIFIC]}
\label{sec:results}
% ============================================================

% Replace this section entirely for each paper.
% Suggested sub-structure:
%   \subsection{Setup / model}
%   \subsection{Main theorem / experiment}
%   \subsection{Numerical validation}

\subsection{Model and setup}
\label{ssec:model}

[Describe the map(s) / system studied.  State parameters used.
For the logistic map: $f(x) = ax(1-x)$, $a \in [1,4]$,
$x \in [0,1]$.]

\subsection{Main result}
\label{ssec:main}

\begin{theorem}[[NAME]]
\label{thm:main}
[State theorem.]
\end{theorem}
%
\begin{proof}
[Proof.]
\end{proof}

\subsection{Numerical verification}
\label{ssec:numerics}

[Describe computation: Python 3.13, NumPy/SciPy, grid sizes,
iteration counts, convergence tolerances.
Reference figures: Fig.~\ref{fig:example}.]

\begin{figure}[htbp]
  \centering
  \includegraphics[width=\columnwidth]{figures/placeholder.pdf}
  \caption{%
    \textbf{[Panel description.]}
    [Detailed caption.  State all parameter values.  Note what
    each panel/colour/linestyle represents.]
  }
  \label{fig:example}
\end{figure}

% ============================================================
\section{Discussion}
\label{sec:discussion}
% ============================================================

[Significance.  Comparison to prior work.  Limitations.
Honest statement of what is proved vs.\ empirically observed.]

% ============================================================
\section{Conclusions}
\label{sec:conclusions}
% ============================================================

[One paragraph per main contribution.  Restate the central
result, the method, and the open questions it raises.]

% ============================================================
\begin{acknowledgments}
The author thanks Howard Lee (in memoriam), Amara Katabarwa,
Eric Suter, and Brandon Campbell for discussions during the
early development of this work.  The term ``Petrification'' was
coined by Eric Suter.
\end{acknowledgments}

% ============================================================
%  APPENDICES
% ============================================================

\appendix

\section{Numerical Methods and Software}
\label{app:numerics}

All computations use Python~3.13 with NumPy~\cite{Numpy2020},
SciPy~\cite{Scipy2020}, and Matplotlib~\cite{Matplotlib2007}.
[Describe: grid sizes, iteration counts, convergence criteria,
random seeds, Dask parallelisation where applicable.]

\section{[ADDITIONAL PROOFS]}
\label{app:proofs}

[Move lengthy derivations here to keep the main text readable.]

% ============================================================
%  BIBLIOGRAPHY
% ============================================================

\bibliography{refs}   % compiles from refs.bib in this directory

\end{document}

% ============================================================
%  ELSARTICLE HEADER (Computer Physics Communications / Elsevier)
%  Uncomment and replace the revtex4-2 header above if needed.
% ============================================================
%
% \documentclass[preprint,12pt,authoryear]{elsarticle}
% \usepackage{lineno}
% \modulolinenumbers[5]
% \journal{Computer Physics Communications}
% \bibliographystyle{elsarticle-harv}
